\section{Appendix A: Reversible Stress from Virtual Displacement}
\label{app:reversible-stress-derivation}

\sourceeqrange{Appendix A Eqs. (A1)--(A4), with Eqs. (2.35), (2.40), (2.47), and (2.48)}

\subsection{Purpose and Scope}

This appendix reconstructs the reversible-stress calculation that the source
places after the main hydrodynamic discussion.  The source Appendix A is short;
the companion expands the typed bookkeeping so that a reader can see which
objects are densities, variations, tensors, and boundary terms.

The scope is narrow.  We derive the virtual-displacement route to the
reversible stress tensor and record its sign convention.  We do not derive the
full hydrodynamic balance laws, do not nondimensionalize Appendix B, and do not
choose between the printed-B3 and dimensional-source branches.

\subsection{Typed Objects and Assumptions}

Let \(\Omega\subset\mathbb{R}^d\) be a smooth container and let
\(x=(x_1,\ldots,x_d)\) be Eulerian coordinates.  The fields used here are
classical and smooth enough for the index manipulations below.  Boundary terms
are displayed before any cancellation assumption is imposed.

\begin{center}
\begin{tabular}{@{}lll@{}}
\toprule
Symbol & Type & Role in this appendix \\
\midrule
\(u:\Omega\to\mathbb{R}^d\) & vector field & virtual displacement \\
\(n:\Omega\to\mathbb{R}_{>0}\) & scalar field & number density \\
\(\delta n\) & scalar variation & density change induced by \(u\) \\
\(\hat e\) & scalar density & gradient-augmented internal-energy density \\
\(\hat S\) & scalar density & gradient-augmented entropy density \\
\(\hat\mu\) & scalar field & generalized chemical potential \\
\(M=M(n,T)\) & coefficient field & square-gradient free-energy coefficient \\
\(\Pi=(\Pi_{ij})\) & tensor field & reversible stress or pressure tensor \\
\(p_1\) & scalar field & diagonal gradient pressure contribution \\
\(\nu\) & boundary vector field & outward unit normal \\
\bottomrule
\end{tabular}
\end{center}

Repeated indices are summed.  The companion writes
\[
  n_i=\partial_i n,\qquad
  u_{i,j}=\partial_j u_i .
\]
All formulas are first-order in \(u\); terms quadratic in the virtual
displacement gradient are omitted.

\subsection{Virtual Displacement and Density Variation}

A virtual displacement sends a material point from \(x\) to
\[
  x'=x+u(x).
\]
The deformation gradient is \(F_{ij}=\partial x_i'/\partial x_j
=\delta_{ij}+\partial_j u_i\).  For a first-order determinant, use
\(\det(I+A)=1+\operatorname{tr}A+O(|A|^2)\).  With
\(A_{ij}=\partial_j u_i\),
\[
  \operatorname{tr}A=A_{ii}=\partial_i u_i=\grad\cdot u,
  \qquad
  \det F=1+\partial_i u_i+O(|\grad u|^2).
\]
so the Jacobian is
\[
  \dd x'=(1+\grad\cdot u)\dd x+O(|\grad u|^2).
\]
For a material volume element, particle conservation \(n'\dd x'=n\dd x\)
gives
\[
  n'=n(1-\grad\cdot u)+O(|\grad u|^2),
\]
and hence the local volume change and material density change recorded in
\eqsrc{A1}:
\begin{equation}
  \delta V=V\,\grad\cdot u,
  \qquad
  \delta_{\mathrm{mat}}n=-n\,\grad\cdot u .\sourceeqbadge{A1}
  \label{eq:app-a-material-density-variation}
\end{equation}
When the same calculation is written over the fixed Eulerian domain, the
induced density variation is
\begin{equation}
  \delta_E n
  =
  -\grad\cdot(nu)
  =
  -u\cdot\grad n-n\,\grad\cdot u .
  \label{eq:app-a-eulerian-density-variation}
\end{equation}
The first equality is the conservation form; the second is the product-rule
form.  It includes both the material compression term \(-n\grad\cdot u\)
and the Eulerian relabeling term \(-u\cdot\grad n\).  The verification
scripts check this identity in components.

\subsection{Coordinate Derivatives and the Gradient Term}

The inverse map has \(x_j=x'_j-u_j(x')+O(|u|^2)\), so the inverse
Jacobian is
\[
  \frac{\partial x_j}{\partial x'_i}
  =\delta_{ij}-\partial_i u_j+O(|\grad u|^2).
\]
The chain rule then gives the operator identity
\[
  \partial_i'=(\partial x_j/\partial x_i')\partial_j.
\]
The coordinate derivative in the deformed coordinate frame therefore has the
first-order form stated in \eqsrc{A3}:
\begin{equation}
  \frac{\partial}{\partial x_i'}
  =
  \frac{\partial}{\partial x_i}
  -
  \frac{\partial u_j}{\partial x_i}
  \frac{\partial}{\partial x_j}
  +O(|\grad u|^2).\sourceeqbadge{A3}
  \label{eq:app-a-derivative-transform}
\end{equation}
Thus the square-gradient density has two kinds of first-order change.  First,
the field changes by \(\delta n\).  Second, the derivative operator changes
with the coordinates.  To first order,
\[
  \partial_i' n'
  =\partial_i n+\partial_i\delta n-(\partial_i u_j)\partial_j n .
\]
Writing out the squared-gradient derivative before contraction gives
\begin{align}
  \frac{1}{2}\delta(\partial_i n\,\partial_i n)
  &=\partial_i n\,\delta(\partial_i n) \\
  &=\partial_i n\left[\partial_i\delta n-(\partial_i u_j)\partial_j n\right].
  \label{eq:app-a-gradient-change-linebyline}
\end{align}
Therefore the derivative of \(|\grad n|^2/2\) is
\begin{equation}
  \delta\!\left(\frac{1}{2}|\grad n|^2\right)
  =
  \grad n\cdot\grad\delta n
  -
  (\partial_i n)(\partial_j n)\,\partial_i u_j .
  \label{eq:app-a-gradient-change}
\end{equation}
Equation \eqref{eq:app-a-gradient-change} is the component mechanism that
later produces the dyadic term \(M\grad n\otimes\grad n\) in the stress
tensor.  The first term is an Euler--Lagrange or boundary-term contribution;
the second term already has the shape of a stress-power contribution.

For the material variation in source Eq. (A1), write
\(\theta_u=\partial_k u_k\) and substitute
\(\delta_{\mathrm{mat}}n=-n\theta_u\).  Its spatial derivative is
\begin{equation}
  \partial_i(\delta_{\mathrm{mat}}n)
  =-(\partial_i n)\theta_u-n\partial_i\theta_u .
  \label{eq:app-a-material-density-gradient}
\end{equation}
The gradient-density variation therefore becomes
\begin{equation}
  \delta\!\left(\frac12|\grad n|^2\right)
  =
  -|\grad n|^2\theta_u
  -n(\partial_i n)\partial_i\theta_u
  -(\partial_i n)(\partial_j n)\partial_i u_j .
  \label{eq:app-a-material-gradient-variation}
\end{equation}
The combination of the entropy and energy gradient coefficients is
\(TC+K=M\).  Hence this contribution to \(T\delta S\) is
\(-M\delta(|\grad n|^2/2)\).  Before any boundary choice, its middle term
satisfies
\begin{align}
  \int_\Omega Mn(\partial_i n)\partial_i\theta_u\,\dd x
  &=-\int_\Omega
    \theta_u\,\partial_i\!\left(Mn\partial_i n\right)\dd x
  \nonumber\\
  &\quad+
  \int_{\partial\Omega}
    Mn(\nu\cdot\grad n)\theta_u\,\dd a .
  \label{eq:app-a-material-gradient-ibp}
\end{align}
Thus the scalar bulk pieces generated by the gradient variation are
\begin{equation}
  \left[
    M|\grad n|^2-\partial_i(Mn\partial_i n)
  \right]\theta_u,
  \label{eq:app-a-material-gradient-scalar-bulk}
\end{equation}
while the anisotropic piece is
\(M(\partial_i n)(\partial_j n)\partial_i u_j\).  The local and coefficient
variations derived in Section II.B, especially
\eqref{eq:iib-muhat} and \eqref{eq:iib-M-relation}, collect the scalar pieces
into the generalized-chemical-potential term used below.  For the interior
stress identification we take \(u\in C_c^2(\Omega;\mathbb{R}^d)\), so the
boundary integral in \eqref{eq:app-a-material-gradient-ibp} vanishes.  For a
wall-touching displacement it must instead be retained and combined with a
specified wall functional; this appendix does not manufacture that global
branch.

\subsection{Energy Variation and Stress Power}

The source uses a nondissipative virtual motion in Appendix A.  In that branch
the internal-energy change of a material element is paired with reversible
mechanical work.  In companion notation, the stress-power convention is
\begin{equation}
  \delta\hat e
  =
  -(\hat e\,\delta_{ij}+\Pi_{ij})\,\partial_i u_j .\sourceeqbadge{A2}
  \label{eq:app-a-energy-work-convention}
\end{equation}
This is the sign convention consistent with the internal-energy balance
\eqsrc{2.40}, where the reversible stress contributes as
\(-\Pi_{ij}\partial_i v_j\).  Because the reversible stress obtained below is
symmetric, the contractions \(\Pi_{ij}\partial_i v_j\) and
\(\Pi_{ij}\partial_j v_i\) are equivalent after dummy-index relabeling.

A useful sign audit is obtained by replacing the virtual displacement gradient
\(\partial_i u_j\) by a velocity gradient \(\partial_i v_j\).  With the convention
above, the reversible contribution to the internal-energy balance is the scalar
\begin{equation}
  -\Pi_{ij}\partial_i v_j,
  \label{eq:app-a-stress-power-sign-audit}
\end{equation}
which has the same sign as the source internal-energy branch used later in
Section~\ref{sec:hydrodynamic-equations}.  The sign convention also fixes the
force-density convention used later in the momentum equation \eqsrc{2.36}:
\begin{equation}
  f_i^{\mathrm{rev}}
  =
  -\partial_j\Pi_{ij}.
  \label{eq:app-a-force-convention}
\end{equation}
The companion records this convention here, but the full momentum equation is
deferred to Section II.D.

\subsection{Entropy Variation and Identification of the Tensor}

Use the local differential from Section II.B in the variables
\((n,\hat e)\),
\[
  \delta\hat S
  =
  \frac{1}{T}\delta\hat e
  -
  \frac{\hat\mu}{T}\delta n
  +\hbox{gradient-coordinate terms}.
\]
Insert the material density variation from \eqsrc{A1}, the work convention
\eqref{eq:app-a-energy-work-convention}, and the square-gradient coordinate
change \eqref{eq:app-a-gradient-change}.  The coefficient collection can be
read term by term:
\begin{align}
  \frac{T\delta S}{V}
  &=T\hat S\,\partial_k u_k
    -\hat e\,\partial_k u_k
    -\Pi_{ij}\partial_i u_j
    +n\hat\mu\,\partial_k u_k
    +M(\partial_i n)(\partial_j n)\partial_i u_j .
  \label{eq:app-a-entropy-variation-uncollected}
\end{align}
The first term is the entropy-density volume contribution, the second and
third terms come from the reversible work convention, the fourth term comes
from the density variation, and the final term comes from the transformed
gradient operator.  The term origins can be checked as follows:
\begin{center}
\begin{tabular}{@{}p{0.35\linewidth}p{0.52\linewidth}@{}}
\toprule
Term & Origin \\
\midrule
\(T\hat S\partial_k u_k\) & volume change of entropy density \\
\(-\hat e\partial_k u_k-\Pi_{ij}\partial_i u_j\) & reversible work convention \\
\(n\hat\mu\partial_k u_k\) & density variation under material compression \\
\(M n_i n_j\partial_i u_j\) & transformed gradient operator \\
\bottomrule
\end{tabular}
\end{center}
To see the index collection, rewrite every divergence contribution as a
coefficient of \(\partial_i u_j\):
\begin{align}
  T\hat S\partial_k u_k&=T\hat S\delta_{ij}\partial_i u_j, \\
  -\hat e\partial_k u_k&=-\hat e\delta_{ij}\partial_i u_j, \\
  n\hat\mu\partial_k u_k&=n\hat\mu\delta_{ij}\partial_i u_j .
  \label{eq:app-a-divergence-to-tensor-ledger}
\end{align}
The dyadic gradient term is already a coefficient of \(\partial_i u_j\).  Thus,
collecting the coefficient of \(\partial_i u_j\), the reversible nondissipative
branch gives
\begin{equation}
  \frac{T\,\delta S}{V}
  =
  \left[
    (T\hat S-\hat e+n\hat\mu)\delta_{ij}
    -
    \Pi_{ij}
    +
    M(\partial_i n)(\partial_j n)
  \right]\partial_i u_j .\sourceeqbadge{A4}
  \label{eq:app-a-entropy-variation-collected}
\end{equation}
For arbitrary virtual displacement gradients in the reversible branch,
\(\delta S=0\).  This means the coefficient of each independent
\(\partial_i u_j\) vanishes.  The independence is local: at an interior point,
a compactly supported virtual displacement can prescribe the first-order
matrix \(\partial_i u_j\) without changing boundary data.  Therefore the
bracketed tensor vanishes, and
\begin{equation}
  \Pi_{ij}
  =
  (T\hat S-\hat e+n\hat\mu)\delta_{ij}
  +
  M(\partial_i n)(\partial_j n).
  \label{eq:app-a-stress-from-entropy}
\end{equation}
This is the companion's derived form of the reversible stress.  It is a local
component identity in the nondissipative virtual-displacement branch; it is
not a proof of the dissipative hydrodynamic system or of any numerical result.

\subsection{Relation to the Section II.D Pressure Tensor}

Section II.D writes the same tensor as a pressure tensor with a diagonal
gradient correction:
\begin{equation}
  \Pi_{ij}
  =
  (p+p_1)\delta_{ij}
  +
  M(\partial_i n)(\partial_j n).\sourceeqbadge{2.47}
  \label{eq:app-a-pi-pressure-form}
\end{equation}
Comparing \eqref{eq:app-a-stress-from-entropy} with
\eqref{eq:app-a-pi-pressure-form} gives the scalar identification
\begin{equation}
  p_1
  =
  n\hat\mu-\hat e+T\hat S-p,\sourceeqbadge{2.48}
  \label{eq:app-a-p1-thermo-form}
\end{equation}
which is the thermodynamic form of the diagonal contribution in source Eq. (2.48).  The
first form of \(p_1\), containing derivatives of \(M\), \(T\), and \(n\), is
part of the Section II.D pressure-tensor branch and is not rescaled here.

\subsection{Diagonal and Off-Diagonal Pieces}

The dyadic gradient term has a simple component split.  For \(i\ne j\),
\begin{equation}
  \Pi_{ij}
  =
  M(\partial_i n)(\partial_j n).
  \label{eq:app-a-off-diagonal}
\end{equation}
For diagonal components,
\begin{equation}
  \Pi_{ii}
  =
  p+p_1+M(\partial_i n)^2
  \qquad\hbox{(no sum over \(i\)).}
  \label{eq:app-a-diagonal}
\end{equation}
The off-diagonal part is symmetric because
\[
  M(\partial_i n)(\partial_j n)
  =
  M(\partial_j n)(\partial_i n).
\]
The diagonal/off-diagonal split used by the verification scripts is therefore
\begin{center}
\begin{tabular}{@{}lll@{}}
\toprule
Component & Diagonal scalar part & Gradient dyad part \\
\midrule
\(i=j\) & \(p+p_1\) & \(M(\partial_i n)^2\) \\
\(i\ne j\) & \(0\) & \(M(\partial_i n)(\partial_j n)\) \\
\bottomrule
\end{tabular}
\end{center}
The verification scripts include an expected nonzero residual when the
off-diagonal gradient term is artificially omitted.  That residual is a
bookkeeping test, not a source branch.

\subsection{Boundary Handling}

Several integrations by parts are available after
\eqref{eq:app-a-eulerian-density-variation}.  For example,
\[
  \int_\Omega \phi\,\grad\cdot(nu)\,\dd x
  =-
  \int_\Omega n u\cdot\grad\phi\,\dd x
  +
  \int_{\partial\Omega}\phi n(u\cdot\nu)\,\dd a .
\]
This appendix does not cancel such terms silently.  Boundary terms vanish only
under an explicit branch such as compactly supported \(u\), fixed boundary
displacement, no-normal displacement, or a separately supplied wall functional.
Section II.C supplied the equilibrium wall density condition; the dynamic wall
and global entropy transfer terms remain later issues.

\subsection{Type Checks and Regularity Assumptions}

\begin{itemize}
\item \(\Pi_{ij}\) and \(p+p_1\) are stress or pressure-density objects.
\item \(M(\partial_i n)(\partial_j n)\) has the same tensor type as a stress
  contribution.
\item \(\Pi_{ij}\partial_i v_j\) is a power density.
\item \(-\partial_j\Pi_{ij}\) is a force density.
\item Boundary traces of \(u\), \(n\), and \(\grad n\) must exist before a
  global virtual-work identity is asserted.
\end{itemize}

\subsection{Verification Checklist}

The companion verification layer for this appendix checks only finite
component identities:

\begin{itemize}
\item the Eulerian density variation
  \(\delta_E n=-\grad\cdot(nu)\);
\item the product-rule expansion of \(\grad\cdot(nu)\);
\item a one-dimensional integration-by-parts sanity check for a gradient
  functional;
\item the component stress-power contraction for
  \((p+p_1)\delta_{ij}+M n_i n_j\);
\item symmetry of the off-diagonal stress term;
\item the expected nonzero residual when the off-diagonal gradient term is
  omitted;
\item a regression row that Appendix-B B3 remains deferred.
\end{itemize}

These checks do not prove thermodynamic consistency, the second law, PDE
well-posedness, the numerical simulations, or any Onuki--Fukagawa global
equivalence claim.

The executable mirrors for the Appendix A derivation are:
\begin{center}
\small
\begin{tabular}{@{}l@{}}
\texttt{scripts/python/appendix\_a\_reversible\_stress\_sympy.py}\\
\texttt{scripts/wolfram/appendix\_a\_reversible\_stress\_checks.wls}\\
\texttt{tests/test\_appendix\_a\_reversible\_stress.py}
\end{tabular}
\end{center}
They check the Eulerian density-variation product rule, the first-order
coordinate derivative transform, the component stress-power contraction, the
off-diagonal symmetry, and the expected nonzero residual when the gradient
dyad is removed.  These finite checks support the local virtual-displacement
calculation only; they do not decide Appendix-B scaling or dynamic boundary
transfer.


\subsection{Source-Procedure Closure for Appendix A}

The completion audit checks Appendix A by source procedure rather than by the
shortness of the printed appendix.  The companion now spells out the kinematic,
coordinate-chain-rule, energy-work, entropy-variation, and tensor-identification
steps before relating the result to the Section II.D pressure tensor.

\begin{center}
\begin{tabular}{@{}p{0.16\linewidth}p{0.30\linewidth}p{0.40\linewidth}@{}}
\toprule
Source item & Procedure type & Companion coverage \\
\midrule
\eqsrc{A1} & displacement kinematics & determinant expansion and Eulerian density variation are displayed \\
\eqsrc{A2} & nondissipative work convention & internal-energy/stress-power sign convention is recorded \\
\eqsrc{A3} and following text & coordinate chain rule and gradient variation & inverse-Jacobian derivative transform and \(|\grad n|^2/2\) variation are expanded \\
\eqsrc{A4} & entropy condition and stress identification & coefficient of \(\partial_i u_j\) is collected and set to zero in the reversible branch \\
\bottomrule
\end{tabular}
\end{center}

This closes the local Appendix A derivation scaffold.  It does not decide any
Appendix-B scaling branch and does not close dynamic boundary-transfer terms.

\subsection*{Tagged verification complement}

\OnukiLinkAppendixA{The tagged Appendix A complement} gives the expanded hand
calculation, finite component checks, and stress-power negative control. It does
not claim a continuum stress derivation or a global boundary map.

\subsection{Open Issues Carried Forward}

\begin{itemize}
\item Appendix-B B3 printed/scaled branch remains unresolved here.
\item Section II.D still has to derive the full hydrodynamic balance-law
  sequence.
\item Wall/global entropy transfer is not closed by the virtual-displacement
  calculation.
\item The unpublished simulation-code branch is not inferred.
\end{itemize}
